\section{Факторное представление для градуировки $D_N$}

Пусть $M$ --- полиномиальное многообразие Фробениуса
размерности $N\geq 5$ с градуировкой $D_N$,
полупростое в общем.

Обозначим
\[
A_0=(T_0M,\circ_0).
\]

Наша цель состоит в построении базиса
\[
\{e_0,e_1,e_1^2,\ldots,e_1^{N-2},e_\tau\}
\]
и последующем представлении семейства касательных
алгебр в виде фактор-алгебры.

Основная трудность заключается в доказательстве
\[
\boxed{e_1^{N-2}\big|_0\neq 0.}
\tag{1}
\]

Здесь $e_1$ --- ненулевой однородный элемент веса
\[
\operatorname{wt}(e_1)=\frac{1}{N-1}.
\]

Условие (1) необходимо и достаточно для того,
чтобы степени
\[
1,e_1,e_1^2,\ldots,e_1^{N-2}
\]
были линейно независимы.

Действительно, все они имеют различные веса, а
ненулевость последней степени влечёт ненулевость
всех предыдущих.

Остающийся базисный элемент можно выбрать в
однородной компоненте веса
\[
\frac{N-2}{2(N-1)}.
\]

При чётном $N$ данная компонента двумерна,
поэтому дополнительный элемент выбирается
линейно независимо от соответствующей степени
$e_1$.

\subsection{Оператор умножения и фильтрация по весам}

Рассмотрим оператор
\[
L_{e_1}:A_0\longrightarrow A_0,
\qquad
L_{e_1}(v)=e_1\circ_0 v.
\]

Градуировка алгебры имеет вид
\[
A_0=\bigoplus_\lambda (A_0)_\lambda.
\]

Поскольку умножение однородно,
\[
L_{e_1}\bigl((A_0)_\lambda\bigr)
\subset
(A_0)_{\lambda+\frac{1}{N-1}}.
\]

Таким образом, оператор $L_{e_1}$ повышает
вес на $\frac{1}{N-1}$.

Для стандартной алгебры Якоби особенности $D_N$
имеется цепочка
\[
1\longmapsto x\longmapsto x^2
\longmapsto\cdots\longmapsto x^{N-2}
\longmapsto 0,
\]
и дополнительный элемент $y$ удовлетворяет
\[
xy=0.
\]

Следовательно, жорданов тип оператора $L_x$
равен
\[
\boxed{(N-1,1).}
\]

Условие (1) эквивалентно существованию у
$L_{e_1}$ жорданова блока размера $N-1$.

Если некоторый структурный коэффициент
обращается в нуль, цепочка степеней
может оборваться и соответствующий жорданов
блок распадается на несколько меньших.

Возникает следующая задача.

\medskip

\noindent
\textbf{Рабочая гипотеза.}
Для полиномиального многообразия Фробениуса
с градуировкой $D_N$, полупростого в общем,
оператор $L_{e_1}(0)$ имеет жорданов блок
размера $N-1$.

\medskip

Это утверждение пока требует доказательства.
Общая полупростота сама по себе не запрещает
вырождения жордановой структуры в специальной
точке.

\section{Пример $D_5$: случай $a=0$}

Рассмотрим подробнее градуировку $D_5$.

Пусть
\[
A_0=T_0M
\]
и выбран однородный базис
\[
e_0,e_1,e_2,e_3,e_\tau
\]
с весами
\[
0,\quad\frac14,\quad\frac12,\quad
\frac34,\quad\frac38.
\]

Зафиксируем нормировку фробениусова
спаривания
\[
\eta(e_0,e_3)=1,
\qquad
\eta(e_1,e_2)=1,
\qquad
\eta(e_\tau,e_\tau)=-4,
\]
причём остальные попарные значения равны нулю.

В точке $0$ градуировка и инвариантность
спаривания дают
\[
e_1\circ_0e_1=ae_2,
\]
\[
e_1\circ_0e_2=e_3,
\]
\[
e_\tau\circ_0e_\tau=-4e_3.
\]

Остальные произведения элементов положительного
веса равны нулю.

Константа $a\in\mathbb C$ не определяется
одной градуировкой и инвариантностью метрики.

Рассмотрим случай
\[
\boxed{a=0.}
\]

Тогда
\[
e_1^2=0,
\qquad
e_1e_2=e_3,
\qquad
e_\tau^2=-4e_3.
\tag{2}
\]

\subsection{Факторное представление в нуле}

Положим
\[
x=e_1,\qquad
y=e_2,\qquad
z=e_\tau.
\]

Из соотношений (2) получаем
\[
x^2=0,\qquad y^2=0,
\]
\[
xz=0,\qquad yz=0,
\]
\[
z^2=-4xy.
\]

Таким образом,
\[
\boxed{
	A_0\simeq
	\frac{\mathbb C[x,y,z]}
	{(x^2,y^2,xz,yz,z^2+4xy)}.
}
\tag{3}
\]

Базис фактор-алгебры имеет вид
\[
\boxed{\{1,x,y,z,xy\}.}
\]

В частности,
\[
\dim_{\mathbb C}A_0=5.
\]

Это градуированная коммутативная
ассоциативная фробениусова алгебра.

Её максимальный идеал
\[
\mathfrak m=(x,y,z)
\]
удовлетворяет
\[
\mathfrak m^2=\langle xy\rangle,
\qquad
\mathfrak m^3=0.
\]

Следовательно,
\[
\dim_{\mathbb C}
\frac{\mathfrak m}{\mathfrak m^2}=3.
\]

Таким образом, для порождения $A_0$
требуются три элемента положительного веса.

\subsection{Неразложимость алгебры}

\begin{prop}
	Алгебра $A_0$, заданная соотношениями (3),
	неразложима в прямое произведение двух
	ненулевых коммутативных алгебр.
\end{prop}

\begin{proof}
	Идеал
	\[
	\mathfrak m=(x,y,z)
	\]
	нильпотентен, причём
	\[
	A_0/\mathfrak m\simeq\mathbb C.
	\]
	
	Следовательно, $A_0$ является локальной
	артиновой алгеброй.
	
	Пусть $p\in A_0$ --- идемпотент:
	\[
	p^2=p.
	\]
	
	В локальной алгебре один из элементов
	$p$ и $1-p$ обратим.
	
	Из равенства
	\[
	p(1-p)=0
	\]
	следует, что $p=0$ или $p=1$.
	
	Таким образом, нетривиальных идемпотентов
	в $A_0$ не существует.
	
	С другой стороны, разложение
	\[
	A_0\simeq B_1\times B_2
	\]
	для ненулевых $B_1,B_2$ порождало бы
	нетривиальный идемпотент $(1,0)$.
	
	Получаем противоречие.
\end{proof}

\begin{remark}
	Следовательно, случай $a=0$ нельзя исключить
	одним лишь требованием локальной
	неразложимости многообразия Фробениуса
	в окрестности начала координат.
	
	Действительно, локальное разложение
	многообразия Фробениуса в прямое произведение
	влекло бы соответствующее разложение
	касательной алгебры в точке $0$.
	
	При этом неразложимость не означает
	полупростоты: рассматриваемая алгебра
	содержит ненулевые нильпотентные элементы.
\end{remark}

\subsection{Отличие от стандартной особенности $D_5$}

Стандартная алгебра Якоби особенности $D_5$
имеет вид
\[
Q_{D_5}=
\frac{\mathbb C[X,Y]}
{(XY,Y^2+4X^3)}.
\]

Для неё
\[
\dim_{\mathbb C}
\frac{\mathfrak m_{D_5}}
{\mathfrak m_{D_5}^2}=2.
\]

В то же время для алгебры (3)
\[
\dim_{\mathbb C}
\frac{\mathfrak m}{\mathfrak m^2}=3.
\]

Поэтому
\[
\boxed{A_0\not\simeq Q_{D_5}.}
\]

Кроме того, в случае $a=0$
\[
\mathfrak m^3=0,
\]
тогда как в стандартной алгебре $D_5$
\[
\mathfrak m_{D_5}^3\neq0.
\]

Жорданов тип оператора умножения на $e_1$
при $a=0$ равен
\[
\boxed{(2,2,1),}
\]
поскольку имеются цепочки
\[
1\longmapsto x\longmapsto0,
\]
\[
y\longmapsto xy\longmapsto0,
\]
\[
z\longmapsto0.
\]

Это отличается от типа $(4,1)$,
соответствующего стандартной алгебре $D_5$.

\section{Использование общей полупростоты}

Обозначим
\[
R=\mathbb C[s_0,\ldots,s_{N-2},\tau]
\]
и рассмотрим семейство полиномиальных
касательных алгебр
\[
A=\Gamma_{\mathrm{poly}}(TM).
\]

Это свободный $R$-модуль ранга $N$,
наделённый коммутативным ассоциативным
умножением.

Общая полупростота означает, что
алгебра над полем рациональных функций
\[
A\otimes_R\operatorname{Frac}(R)
\]
является конечной этальной алгеброй.

После перехода к алгебраическому замыканию
поля $\operatorname{Frac}(R)$ она изоморфна
\[
\overline{\operatorname{Frac}(R)}^{\,N}.
\]

Рассмотрим относительный спектр
\[
\mathcal X=\operatorname{Spec}_R A
\]
и естественный конечный плоский морфизм
\[
\pi:\mathcal X\longrightarrow
\operatorname{Spec}R.
\]

В общей полупростой точке слой $\pi$
состоит из $N$ различных геометрических точек.

Вопрос состоит в том, может ли градуированная
алгебра, у которой цепочка степеней $e_1$
обрывается, возникать в качестве специального
слоя такого семейства при выполнении
условий многообразия Фробениуса.

В частности, для $D_5$ требуется исследовать,
возможно ли
\[
e_1^2\big|_0=0
\]
при общей полупростоте семейства.

Само по себе существование нильпотентных
элементов в специальном слое не противоречит
общей полупростоте.

Таким образом, необходимо использовать
дополнительные ограничения, связывающие
умножение в различных точках, в частности
потенциальность и ассоциативность.

\subsection{Следовая форма и каустика}

Для $v\in A$ обозначим через $L_v$
оператор умножения на $v$.

Рассмотрим следовую форму
\[
B(v,w)=\operatorname{Tr}(L_vL_w).
\]

В полиномиальном базисе $e_i$
определим
\[
\Delta(s)=
\det\bigl(
\operatorname{Tr}(L_{e_i}L_{e_j})
\bigr).
\]

Поскольку основное поле имеет
характеристику нуль,
\[
\boxed{
	M\text{ полупросто в общем}
	\quad\Longleftrightarrow\quad
	\Delta\not\equiv0.
}
\]

Поэтому один из возможных путей состоит
в доказательстве импликации
\[
e_1^{k+1}\big|_0=0
\quad\Longrightarrow\quad
\Delta\equiv0.
\]

Однако эта импликация пока не установлена
и не следует из одной градуировки.

\section{Возможная схема доказательства}

Предположим, что
\[
e_1^k\big|_0\neq0,
\qquad
e_1^{k+1}\big|_0=0
\]
для некоторого $k<N-2$.

Можно попытаться исследовать последствия
первого обрыва цепочки степеней.

\medskip

\noindent
\textbf{Первый этап.}
Классифицировать допустимые жордановы
типы $L_{e_1}(0)$ с учётом градуировки
и инвариантного фробениусова спаривания.

\medskip

\noindent
\textbf{Второй этап.}
Для первого обрыва цепочки определить,
какие структурные константы обязательно
обращаются в нуль вследствие ассоциативности.

\medskip

\noindent
\textbf{Третий этап.}
Исследовать, препятствует ли этот обрыв
существованию полупростого общего слоя
при условии потенциальности умножения.

В частности, можно искать ненулевой
глобальный идеал, все элементы которого
нильпотентны в каждой точке.

Существование такого идеала дало бы
\[
\Delta\equiv0.
\]

Однако его существование ещё необходимо
доказать; оно не следует автоматически
из нильпотентности элементов в $A_0$.

\medskip

\noindent
\textbf{Четвёртый этап.}
Исключив все возможные обрывы цепочки,
получить
\[
e_1^{N-2}\big|_0\neq0.
\]

После этого построить базис
\[
\{e_0,e_1,e_1^2,\ldots,e_1^{N-2},e_\tau\}
\]
на всём многообразии.

\section{Переход к факторному представлению}

Предположим, что удалось доказать
\[
e_1^{N-2}\big|_0\neq0.
\]

Выберем однородные полиномиальные поля
$X,Y$ так, чтобы их специализации в нуле
совпадали с построенными образующими.

Тогда элементы
\[
1,X,\ldots,X^{N-2},Y
\]
образуют базис в точке $0$.

Матрица перехода от исходного
однородного базиса имеет определитель
веса ноль.

Поскольку все координаты базы имеют
положительные веса, этот определитель
является константой.

Его значение в нуле отлично от нуля,
следовательно, он ненулевой всюду.

Поэтому
\[
\boxed{
	A=
	R\cdot 1\oplus
	\bigoplus_{k=1}^{N-2}R\cdot X^k
	\oplus R\cdot Y.
}
\]

В частности, существует соотношение
\[
\boxed{
	X^{N-1}
	=
	\sum_{k=0}^{N-2}
	\alpha_k(s,\tau)X^k
	+\beta(s,\tau)Y.
}
\tag{4}
\]

Аналогично, можно разложить
\[
XY,\qquad Y^2
\]
по построенному базису.

Это даёт соотношения редукции
и возможность представить $A$
в виде фактор-алгебры
\[
A\simeq R[X,Y]/I.
\]

Таким образом, задача получения
факторного представления сводится
к доказательству ненулевости
$e_1^{N-2}$ в начале координат.

\section{Возможность ослабления цели}

Для последующего применения условия
Хертлинга--Манина существенным может
оказаться не всё факторное представление,
а только соотношение редукции (4).

Поэтому следует отдельно проверить,
достаточно ли информации о минимальном
многочлене оператора $L_{e_1}$ вместо
построения полного базиса из степеней.

Если удастся ослабить необходимую
структуру умножения, можно избежать
доказательства максимальной длины
жордановой цепочки.

\section{Заключение}

Для $D_5$ случай $a=0$ приводит к
неразложимой локальной фробениусовой
алгебре
\[
\frac{\mathbb C[x,y,z]}
{(x^2,y^2,xz,yz,z^2+4xy)},
\]
которая не изоморфна стандартной
алгебре Якоби $D_5$.

Следовательно, градуировка,
ассоциативность и невырожденное
инвариантное спаривание сами по себе
не обеспечивают существование
требуемого базиса из степеней.

Основная задача состоит в том, чтобы
выяснить, исключают ли потенциальность
и общая полупростота подобные
вырождения для произвольного $D_N$.

Наиболее желателен аргумент,
основанный на общей структуре
оператора $L_{e_1}$, а не на
явном решении полной системы WDVV
для каждой размерности.